%%
% 引言或背景
% 引言是论文正文的开端,应包括毕业论文选题的背景、目的和意义;对国内外研究现状和相关领域中已有的研究成果的简要评述;介绍本项研究工作研究设想、研究方法或实验设计、理论依据或实验基础;涉及范围和预期结果等。要求言简意赅,注意不要与摘要雷同或成为摘要的注解。
% modifier: 黄俊杰(huangjj27, 349373001dc@gmail.com)
% update date: 2017-04-15
%%

\chapter{绪论}
%定义，过去的研究和现在的研究，意义，与图像分割的不同,going deeper
\label{cha:introduction}
\section{课题背景及研究意义}
\label{sec:background}
% What is the problem
% why is it interesting and important
% Why is it hards, why do naive approaches fails
% why hasn't it been solved before
% what are the key components of my approach and results, also include any specific limitations，do not repeat the abstract
%contribution
\subsection{课题背景}
随着汽车工业与智能交通技术的飞速发展，高级驾驶辅助系统（ADAS）已成为提升行车安全、缓解交通压力的重要技术方向。根据世界卫生组织发布的道路交通伤害事实报告，全球每年约有135万人死于道路交通事故，因而围绕环境感知与辅助决策的安全技术研究具有突出的现实意义~\cite{who2023}。车道线检测作为ADAS环境感知模块的关键组成部分，其核心任务是通过车载摄像头等传感器获取道路图像，精准识别车道边界并输出位置信息，为车道保持辅助、自动变道等功能提供核心数据支撑。

在理想天气条件下，传统车道线检测方法，如边缘检测、霍夫变换及主流深度学习方法已能实现较高的检测精度~\cite{neven2018}。但在实际行车场景中，恶劣天气会严重干扰检测性能，其中雨天是最为常见的挑战场景之一。雨天环境下，雨滴在摄像头镜头上形成遮挡、路面水膜产生镜面反射、车道线与路面的颜色对比度显著降低，同时积水区域、路面坑洼等因素会进一步破坏车道线的完整性与连续性，导致现有检测算法出现漏检、误检等问题。对于 ADAS 而言，一旦车道线识别环节失稳，后续的车道保持、偏离预警和轨迹约束都会受到连锁影响，因此雨天场景下的鲁棒感知能力具有直接的安全意义。

当前，国内外虽已围绕复杂环境下的车道线检测开展了较多研究，但面向雨天条件的鲁棒检测仍存在明显不足：部分算法主要针对轻度降雨场景优化，在中大雨等条件下性能下降较快；一些深度学习模型虽然提升了检测精度，却仍面临计算复杂度较高、实时性不足的问题，难以稳定满足 ADAS 对 30 FPS 以上处理速度的要求；与此同时，覆盖不同降雨强度与道路类型的标准化雨天车道线数据集仍相对缺乏，算法泛化能力因此受到限制。基于上述原因，开展面向 ADAS 的雨天车道线鲁棒检测研究，具有明确的现实紧迫性。

\subsection{研究意义}
\subsubsection{理论意义}
本课题围绕雨天环境下车道线视觉特征退化这一核心问题，分析雨滴、水膜等干扰因素的作用机制，并在此基础上探索面向检测任务的数据组织与特征增强方法。相关研究有助于深化对低信噪比、低对比度条件下细长结构目标检测问题的认识，也可为复杂环境视觉感知中的特征提取与融合策略提供参考。此外，本课题构建的雨天车道线专用数据集，可在一定程度上弥补现有数据集对恶劣天气覆盖不足的问题，为后续研究提供较为直接的数据支撑。

\subsubsection{现实意义}
本课题形成的雨天鲁棒检测方法可直接服务于高级驾驶辅助系统，有助于提升雨天条件下车道线检测的准确性与实时性，从而降低因车道识别失效带来的潜在风险。进一步看，相关成果也可为自动驾驶环境感知模块在复杂天气下的稳定运行提供技术参考，并可延伸至智能交通监测等场景，为雨天道路运行状态分析与交通管理提供辅助数据，因而具有一定的工程应用价值。

\section{雨天车道线检测国内外研究现状}
\label{sec:related_work}
\subsection{国内研究现状}
从研究脉络来看，车道线检测已由基于边缘、几何约束和模型拟合的传统方法，逐步演进到基于深度学习的分割式、锚点式和参数回归式方法。相关综述指出，道路结构先验、场景复杂度以及实时性约束始终是该方向的三大核心问题~\cite{barhillel2014survey,cao2019lane}。国内在车道线检测领域的研究发展迅速，尤其在复杂环境优化方面取得了显著进展。针对复杂场景下的车道检测问题，Zheng等人提出了RESA（Recurrent Feature-Shift Aggregator）方法，通过沿特征图的多方向递归信息传播增强车道线的结构连续性，从而提升了算法对弯道与遮挡场景的适应性~\cite{zheng2021}。

多传感器融合技术是提升环境感知鲁棒性的关键方向，相关研究为全天候车道线检测提供了重要支撑。研究表明，单一传感器在恶劣天气下存在性能局限，例如摄像头在雨天易受雨滴遮挡、水膜反射影响，而毫米波雷达和激光雷达则能提供稳定的距离、速度等数据。通过融合相机、激光雷达、毫米波雷达等多源传感器数据，可实现优势互补，提升复杂环境下感知的准确性和可靠性。相关技术已在国产ADAS系统中得到应用，有效推动了其环境适应性的提升~\cite{zhu2024}。与此同时，针对复杂道路和动态场景，国内外也在持续探索更强的结构建模与上下文聚合机制，为雨天条件下的稳健检测提供方法储备~\cite{li2022,zheng2022clrnet,hou2019sad,qin2020ultrafast}。

针对恶劣天气下的检测难题，国内外学者在算法设计与工程应用上均开展了工作。Li等人提出了Line-CNN方法，通过引入Line Proposal Unit对细长车道线结构进行更有效的建模，提升了端到端车道线检测的精度与实时性~\cite{li2022}。在此基础上，VPGNet通过引入消失点引导的多任务学习框架，较早关注了雨天、夜间等复杂场景中的车道线与道路标识联合识别问题~\cite{lee2017vpgnet}；EL-GAN则从结构一致性出发，通过嵌入损失约束改善分割结果的形状质量，减轻传统语义分割后处理对启发式规则的依赖~\cite{ghafoorian2018elgan}。

在算法层面，有研究专注于利用深度学习提升模型在复杂光照及恶劣天气下的性能。例如，张荣等人针对低光照场景设计了鲁棒的实时车道线检测方法，通过优化网络结构强化模型对低信噪比环境下车道线特征的提取能力，其核心思路对雨天低对比度场景的算法设计具有重要参考价值~\cite{zhang2024}。还有研究探索了结合生成对抗网络（GAN）提升车道线检测模型的鲁棒性，如潘玉恒等人提出引入自注意力生成对抗网络（SAGAN）生成多种环境下的车道线图像以扩充数据集，有效提升了模型在视觉遮挡和多变照明条件下的检测性能，同时避免了训练过拟合，为雨天等恶劣天气场景下的模型泛化能力提升提供了有效方案~\cite{pan2024}。从检测范式上看，轻量化自蒸馏模型ENet-SAD、行列混合锚点分类模型Ultra Fast Lane Detection及其后续改进版本，分别从监督信号增强、全局上下文建模与速度优化角度推动了工程可部署性提升~\cite{hou2019sad,qin2020ultrafast,qin2024ultrafastv2}。

\subsection{国外研究现状}
国外关于车道线检测的研究起步较早，在复杂环境适应性优化方面积累了较多成果。早期研究主要聚焦于传统图像处理方法的改进。随着深度学习技术的发展，基于卷积神经网络（CNN）的车道线检测方法成为主流。Pan等人提出的SCNN通过在特征图空间维度上传播信息，显著增强了车道线这类细长结构的连续建模能力，为复杂环境下的车道检测提供了有效思路~\cite{pan2018spatial}。随后，端到端实例分割方法、注意力引导方法和跨层细化方法相继出现，分别代表了基于分割、基于局部几何约束以及基于多尺度融合的三条典型技术路线~\cite{neven2018,tabellion2021laneatt,zheng2022clrnet}。

针对恶劣天气下的感知挑战，Porav等人研究了图像去雨恢复问题，证明在雨纹和雨滴干扰显著的场景中，针对性图像恢复能够改善下游视觉任务的输入质量~\cite{porav2019}。Fu等人、Yang等人以及Li等人提出的深度去雨方法，则分别从残差学习、联合雨纹检测和多阶段上下文聚合等角度提升了单幅图像去雨质量，为雨天车道检测提供了较成熟的前端增强思路~\cite{fu2017ddn,yang2017jointderain,li2018rescan}。这类研究为理解并缓解雨天视觉退化问题提供了有价值的参考。

在应对雨天等特定干扰方面，Zhang等人提出了一种基于物理启发图像合成的去雨方法，能够有效恢复被雨滴和雨纹遮挡的图像内容~\cite{zhang2018}。近年的三维车道线检测研究进一步表明，仅依赖二维图像中局部纹理并不足以覆盖起伏路面、坡道或遮挡场景，需要同时引入透视变换、拓扑关系和空间几何先验~\cite{garnett2019threedlanenet,chen2022persformer}。此类图像复原与空间建模技术的进步，为提升雨天车道线检测的前端输入质量和整体场景理解能力提供了技术支持。

此外，条件卷积与动态参数生成等新型检测架构也在车道线检测中展现出潜力。Liu等人提出CondLaneNet，将车道线检测建模为条件卷积驱动的自顶向下实例推理过程，在多个基准数据集上取得了优异性能~\cite{liu2021}。这类方法对结构先验与全局上下文的综合利用，有助于应对雨天局部遮挡导致的特征不连续问题。综合来看，当前研究正在从“单纯提高静态精度”转向“兼顾复杂天气鲁棒性、几何一致性与嵌入式实时性”的综合优化，这也构成了本课题选取轻量化增强与检测协同路线的直接依据~\cite{qin2024ultrafastv2,hou2019sad,sandler2018mobilenetv2,jacob2018quantization}。

\subsection{研究现状述评}
综合国内外研究进展可以看出，车道线检测已由传统图像处理方法逐步转向深度学习方法，复杂环境下的适应能力也成为持续关注的重点。现有工作大致沿三条路径展开：一是通过模型结构设计提升对细长目标及上下文关系的建模能力；二是通过数据集构建、数据合成与增强扩大复杂场景覆盖范围；三是通过图像复原、特征增强与轻量化设计兼顾恶劣天气鲁棒性和部署效率。

然而，面向 ADAS 应用的雨天车道线检测仍存在几项关键问题：其一，现有研究对雨滴遮挡、水膜反射和低对比度退化等干扰的讨论，多停留在现象层面，与车道线特征提取过程之间的对应关系仍不够清晰；其二，图像增强与检测模块往往被分离设计，前端复原结果与后端检测目标之间容易出现偏差；其三，雨天专用数据基础相对薄弱，尤其缺少能够同时服务于训练、验证和部署评估的成体系样本组织；其四，不少在复杂场景中表现较好的方法，仍难兼顾车载平台对实时性、稳定性与部署成本的要求。

围绕上述问题，本文按照“问题分析 - 数据构建 - 方法设计 - 工程验证”的思路展开研究：先分析雨天条件下车道线视觉退化的主要表现及其对检测任务的影响，再构建并整理面向雨天场景的 RainLane 数据集，随后在 CLRNet 基线之上设计轻量化抗干扰模块，最后结合实验结果与部署链路验证方案的可行性。按照这一思路展开，有助于使各章节围绕同一研究目标逐步推进。

\section{主要研究内容与结构安排}
\label{sec:arrangement}

本论文共分为六章，各章内容安排如下。

第一章 绪论。介绍课题的研究背景与研究意义，综述国内外在车道线检测及恶劣天气环境感知等方面的研究现状，并说明本文的研究内容与整体结构。

第二章 相关理论与技术基础。系统介绍车道线检测方法、恶劣天气图像增强技术以及深度学习轻量化与边缘端部署技术，为后续数据集构建、算法设计与工程验证提供理论支撑。

第三章 雨天车道线数据集构建与增强。详细说明数据采集方案、标注规范与质量控制流程，分析雨天干扰因素对车道线视觉特征的影响，并介绍数据增强流水线的设计与实现。

第四章 面向雨天的车道线检测算法设计。分析 CLRNet 基线模型的结构特点及其在雨天场景下的失效原因，设计并集成“频域门控 + 方向注意力”抗干扰模块，并说明其与现有训练目标和特征提取流程的融合方式。

第五章 模型部署与实验验证。介绍模型转换与推理优化流程，说明在 NVIDIA Orin 平台上的部署实现，并给出整体性能、分场景性能、可视化结果以及边缘端部署测试结果。

第六章 总结与展望。总结本文的主要研究成果与创新点，指出当前工作的不足，并对后续研究方向进行展望。

\begin{figure}[htbp]
	\centering
	\hspace{-3.8cm}
	\includegraphics[width=1.1\textwidth]{image/chap01/组织架构图.png}
	\caption{本文研究组织架构图}
	\label{fig:construction}
\end{figure}
